%% kapa: exam=1 round=30 year=2019 period=2 form=A subject=accounting subject_ko=회계학 seq=1 number=41
%% source: Q-Net 감정평가사 기출문제 (cst00309) artlSeq=5205205 file=extracted/1008/5205205/2019년 제30회 감정평가사 1차시험 문제지/2019년 제30회 감정평가사 1차시험 문제지/2019년도 제30회 감정평가사 1차 2교시 A형.hwp
%% conversion: kordoc 4.24.51 -> tools/gen1.py (auto); review=proofread-sonnet; flags=none
\begin{question}{41}
투자부동산에 관한 설명으로 옳지 않은 것은?
\begin{choices}
  \item 미래에 투자부동산으로 사용하기 위하여 건설 또는 개발 중인 부동산은 투자부동산에 해당한다.
  \item 소유 투자부동산은 최초 인식시점에 원가로 측정하며, 거래원가는 최초 측정치에 포함한다.
  \item 통상적인 영업과정에서 판매하기 위한 부동산이나 이를 위하여 건설 또는 개발 중인 부동산은 투자부동산에 해당하지 않는다.
  \item 투자부동산을 개발하지 않고 처분하기로 결정하는 경우에는 재고자산으로 재분류한다.
  \item 투자부동산에 대하여 공정가치모형을 선택한 경우, 투자부동산의 공정가치 변동으로 발생하는 손익은 발생한 기간의 당기손익에 반영한다.
\end{choices}
\end{question}
\answer{4}
